% ════════════════════════════════════════════════════════════════
\subsection{Matroidy}
% ════════════════════════════════════════════════════════════════

\begin{definition}
	\uemph{Matroid} to para $M = (S, \mathcal{A})$ taka, że:
	\begin{enumerate}
		\item $S$ -- niepusty zbiór skończony,
		\item $\mathcal{A} \subseteq 2^S$ i $\mathcal{A} \neq \emptyset$ -- czyli $\mathcal{A}$ to niepusta rodzina podzbiorów $S$,
		\item \label{matroid:dziedziczność} $\forall A, B \subseteq S \quad A \in \mathcal{A} \wedge B \subseteq A \implies B \in \mathcal{A}$,
		\item \label{matroid:wymiana} własność wymiany:
		      \begin{flalign*}
			       & \forall A, B \subseteq S \quad A, B \in \mathcal{A} \wedge |B| > |A| \implies \text{istnieje } x \in B \setminus A \text{ taki, że } A \cup \{x\} \in \mathcal{A}. &  &  &  &
		      \end{flalign*}
	\end{enumerate}
	$\mathcal{A}$ -- rodzina zbiorów niezależnych matroidu $M$, elementy $\mathcal{A}$ -- zbiory niezależne matroidu $M$.
\end{definition}

\begin{example} \leavevmode
	\begin{enumerate}
		\item $A$ -- macierz, $S$ -- zbiór wierszy $A$, $\mathcal{A} = \{R \subseteq S : R \text{ tworzy zbiór liniowo niezależny}\}$.\\
		      Wtedy $(S, \mathcal{A})$ to matroid.
		\item $G = (V, E)$ -- graf, $\mathcal{A} = \{F \subseteq E : G[F] \text{ jest acykliczny}\}$.\\
		      Wtedy $(E, \mathcal{A})$ to matroid zwany matroidem grafowym.
	\end{enumerate}
	Z definicji $\emptyset \in \mathcal{A}$.
\end{example}

\begin{lemma} \label{lemma:matroid:liczność}
	Niech $M = (S, \mathcal{A})$. Wszystkie maksymalne zbiory niezależne $M$ mają tę samą liczność.
\end{lemma}
\begin{proof}
	Niech $A, B \in \mathcal{A}$ będą maksymalnymi zbiorami niezależnymi $M$. Przypuśćmy, że $|B| > |A|$.\\
	Z własności \ref{matroid:wymiana}. istnieje $x \in B \setminus A$ taki, że $A \cup \{x\} \in \mathcal{A}$ -- sprzeczność z maksymalnością $A$.
\end{proof}

\begin{problem}
	$M = (S, \mathcal{A})$ -- matroid, $w: S \to \Rp$ -- funkcja wagi.\\
	Dla $A \subseteq S$ niech $w(A) = \sum_{x \in A} w(x)$ -- waga $A$.\\
	Znaleźć zbiór niezależny o maksymalnej wadze w $M$.\\
	$S = \{x_1, x_2, \ldots, x_n\}$.
\end{problem}

\begin{alg}{Algorytm Greedy}{alg:greedy}
	\FunctionNN{Greedy}{$M, w$}
	\State $n \Assign |S|$
	\State $A \Assign \emptyset$ \Comment{$A$ -- rozwiązanie}
	\State posortuj $S$ w porządku nierosnącym względem wagi, tak żeby $w(x_1) \geq w(x_2) \geq \ldots \geq w(x_n)$
	\For{$i \Assign 1$ \To $n$}
	\If{$A \cup \{x_i\} \in \mathcal{A}$} \label{line:greedy-niezaleznosc}
	\State $A \Assign A \cup \{x_i\}$
	\EndIf
	\EndFor
	\State \Return $A$
	\EndFunction
\end{alg}

\begin{theorem}[Rado, Edmonds] \label{thm:rado-edmonds}
	Algorytm Greedy znajduje zbiór niezależny o maksymalnej wadze.
\end{theorem}
\begin{proof}
	Niech $A = \{x_{i_1}, x_{i_2}, \ldots, x_{i_k}\}$ będzie zbiorem zwróconym przez algorytm i $w(x_{i_1}) \geq w(x_{i_2}) \geq \ldots \geq w(x_{i_k})$.\\
	Najpierw pokażemy, że $A$ jest maksymalnym zbiorem niezależnym.\\
	Przypuśćmy, że istnieje $x_j \in S \setminus A$ taki, że $A \cup \{x_j\} \in \mathcal{A}$.
	Wtedy $j < i_k$, bo w przeciwnym wypadku algorytm dodałby $x_j$ do $A$ jako $(k+1)$-szy element.\\
	Niech $l$ będzie najmniejszą liczbą taką, że $j < i_l$.
	Wtedy $w(x_{i_1}) \geq w(x_{i_2}) \geq \ldots \geq w(x_{i_{l-1}}) \geq w(x_j) \geq w(x_{i_l})$.\\
	$A \cup \{x_j\} \in \mathcal{A}$, więc z własności \ref{matroid:dziedziczność}. $\{x_{i_1}, x_{i_2}, \ldots, x_{i_{l-1}}, x_j\} \subseteq A \cup \{x_j\}$ jest zbiorem niezależnym.\\
	Ale wtedy algorytm dodałby $x_j$ zamiast $x_{i_l}$ jako $l$-ty element dodany do $A$ -- sprzeczność.\\
	Zatem $A$ jest maksymalnym zbiorem niezależnym.

	Rozważmy zbiór niezależny $B = \{y_1, y_2, \ldots, y_m\} \in \mathcal{A}$ i załóżmy, że $w(y_1) \geq w(y_2) \geq \ldots \geq w(y_m)$.\\
	Mamy $m \leq k$, bo $A$ jest maksymalnym zbiorem niezależnym (Lemat \ref{lemma:matroid:liczność}).\\
	Pokażemy, że $w(B) \leq w(A)$.\\
	Pokażemy, że dla każdego $j \leq m \quad w(y_j) \leq w(x_{i_j})$.\\
	Przypuśćmy, że $w(y_j) > w(x_{i_j})$.\\
	Rozważmy zbiory niezależne $A_{j-1} = \{x_{i_1}, \ldots, x_{i_{j-1}}\} \subseteq A$, $B_j = \{y_1, \ldots, y_j\} \subseteq B$.\\
	$|B_j| = j > |A_{j-1}| = j-1$, więc z własności \ref{matroid:wymiana}. istnieje $y_s \in B_j \setminus A_{j-1}$ taki, że $A_{j-1} \cup \{y_s\} \in \mathcal{A}$.\\
	Mamy $w(y_s) \geq w(y_j) > w(x_{i_j})$, więc istnieje $t \leq j$ takie, że $w(x_{i_1}) \geq w(x_{i_2}) \geq \ldots \geq w(x_{i_{t-1}}) \geq w(y_s) > w(x_{i_t})$.\\
	Z własności \ref{matroid:dziedziczność}. $\{x_{i_1}, x_{i_2}, \ldots, x_{i_{t-1}}, y_s\} \subseteq A_{j-1} \cup \{y_s\}$ jest zbiorem niezależnym.\\
	Ale wtedy algorytm dodałby $y_s$ zamiast $x_{i_t}$ jako $t$-ty element dodany do $A$ -- sprzeczność.\\
	Zatem
	$\begin{rcases}
			w(y_1) \leq w(x_{i_1}) \\
			w(y_2) \leq w(x_{i_2}) \\
			\vdots                 \\
			w(y_m) \leq w(x_{i_m})
		\end{rcases} \implies w(B) = w(y_1) + \ldots + w(y_m) \leq w(x_{i_1}) + \ldots + w(x_{i_m}) \leq w(A)$.
\end{proof}

\paragraph{Złożoność czasowa algorytmu Greedy}
$n = |S|$
\begin{itemize}[nosep]
	\item linia 1: $\BO(1)$
	\item linia 2: $\BO(1)$
	\item linia 3: $\BO(n \log n)$
	\item liczba iteracji pętli z linii 4 wynosi $n$
	\item linia 5: $f(n)$ -- czas sprawdzenia warunku z linii 5
	\item linia 6: $\BO(1)$
	\item linia 7: $\BO(1)$
\end{itemize}

% ════════════════════════════════════════════════════════════════
\subsection{Problem szeregowania zadań}
% ════════════════════════════════════════════════════════════════

\begin{itemize}[nosep]
	\item $S = \{z_1, z_2, \ldots, z_n\}$ -- zbiór zadań o jednostkowym czasie wykonania
	\item $d_1, d_2, \ldots, d_n$ -- deadliny, $d_i \in \{1, \ldots, n\}$ dla $i \in [n]$
	\item $w_1, w_2, \ldots, w_n$ -- kary, $w_i \geq 0$ dla $i \in [n]$
\end{itemize}
Płacimy karę $w_i$ za zadanie $z_i$, jeżeli nie jest skończone w momencie $d_i$.\\
\begin{problem}[szeregowania zadań z karami]
	Znaleźć kolejność wykonania zadań minimalizującą sumę kar.
\end{problem}

Załóżmy, że mamy pewien harmonogram. Zadanie $z_i$ nazywamy \uemph{terminowym} w tym harmonogramie, jeśli jest zakończone do chwili $d_i$, w przeciwnym razie nazywamy je \uemph{spóźnionym}.\\
Nie płacimy kar za zadania terminowe, płacimy kary za zadania spóźnione.

\begin{lemma}
	Każdy harmonogram można przetransformować na harmonogram z taką samą sumą kar, w którym wszystkie zadania terminowe poprzedzają wszystkie zadania spóźnione oraz zadania terminowe są posortowane w kolejności niemalejącej względem deadlinów.
\end{lemma}
\begin{proof}~
	\begin{enumerate}
		\item Niech $z_i$ będzie zadaniem spóźnionym ukończonym w chwili $l_i > d_i$ i niech $z_j$ będzie zadaniem terminowym ukończonym w chwili $l_j \leq d_j$ i niech $z_i$ będzie przed $z_j$ w harmonogramie, tzn. $l_i < l_j$.\\
		      Zamieniamy miejscami $z_i$ i $z_j$. Po zamianie:
		      \begin{itemize}[nosep]
			      \item $z_i$ jest ukończone w chwili $l_j > l_i > d_i$, więc nadal jest spóźnione.
			      \item $z_j$ jest ukończone w chwili $l_i < l_j \leq d_j$, więc nadal jest terminowe.
		      \end{itemize}
		      Po pewnej liczbie takich zamian otrzymujemy harmonogram o tej samej sumie kar, w którym wszystkie zadania terminowe poprzedzają wszystkie zadania spóźnione.

		\item Niech $z_i, z_j$ będą zadaniami terminowymi takimi, że $z_i$ jest ukończone w chwili $l_i \leq d_i$, $z_j$ jest ukończone w chwili $l_j \leq d_j$, $l_i < l_j$, ale $d_i > d_j$.\\
		      Zamieniamy miejscami $z_i$ i $z_j$. Po zamianie:
		      \begin{itemize}[nosep]
			      \item $z_i$ jest ukończone w chwili $l_j \leq d_j < d_i$, więc $z_i$ jest nadal terminowe.
			      \item $z_j$ jest ukończone w chwili $l_i < l_j \leq d_j$, więc $z_j$ jest nadal terminowe.
		      \end{itemize}
		      Po pewnej liczbie takich zamian otrzymujemy harmonogram spełniający warunki lematu.
	\end{enumerate}
\end{proof}
